% !PRINTED 243
\weberpage{243}
Um eine Grundlage für diese Theorie zu gewinnen, betrachten wir eine ganze Function einer veränderlichen Grösse $t$

\weq{1}{f(t) = a_0 t^n + a_1 t^{n-1} + \cdots + a_{n-1} t + a_n,}

deren Co\"efficienten ganze Zahlen sein sollen. Wir wählen ausserdem eine Primzahl $p$ als Modulus, und nehmen an, $a_0$ sei durch $p$ nicht theilbar.

Zwei ganze Functionen von $t$, in denen entsprechende Co\"efficienten nach dem Modul $p$ congruent sind, sollen selbst \gesperrt{nach dem Modul $p$ congruent} genannt werden.

Die Function $f(t)$ heisst \gesperrt{nach dem Modul $p$ reducibel}, wenn es zwei ganze ganzzahlige Functionen $\varphi(t)$, $\psi(t)$ giebt, in deren jeder wenigstens ein von $t$ abhängiges Glied einen durch $p$ nicht theilbaren Co\"efficienten hat, so dass

\weq{2}{f(t) \equiv \varphi(t)\,\psi(t)\ (\mathrm{mod}\ p)}

ist. In $\varphi(t)$ und $\psi(t)$ kann man alle Glieder, deren Co\"efficienten durch $p$ theilbar sind, weglassen, und dann muss der Grad von $\varphi(t)\,\psi(t)$ mit dem Grade von $f(t)$ übereinstimmen. Es müssen also die Grade von $\varphi(t)$ und $\psi(t)$ niedriger als $n$ sein. Giebt es keine solche Functionen $\varphi(t)$, $\psi(t)$, so heisst $f(t)$ \gesperrt{nach dem Modul $p$ irreducibel}.

Eine nach dem Modul $p$ irreducible Function ist gewiss immer im Körper der rationalen Zahlen absolut irreducibel. Das Umgekehrte ist aber nicht nothwendig.

So ist z.\,B. die Function $2^{\mathrm{ten}}$ Grades $t^2 - R$ reducibel nach $p$, wenn $R$ quadratischer Rest von $p$ ist; denn ist $a^2 \equiv R\ (\mathrm{mod}\ p)$, so ist

\[
t^2 - R \equiv (t-a)(t+a)\ (\mathrm{mod}\ p).
\]

Dagegen ist $t^2 - N$, wenn $N$ Nichtrest von $p$ ist, irreducibel, weil sonst $t^2 - N$ für ein rationales $t$ durch $p$ theilbar werden müsste.

Ein anderes Beispiel bieten die Functionen

\[
t^2 + t + 1, \quad t^3 + t + 1,
\]

die für den Modul $2$ irreducibel sind. Denn wären sie reducibel, so müsste einer der Factoren linear sein, und es müsste eine ganze Zahl geben, die, für $t$ eingesetzt, diese Functionen zu geraden Zahlen macht. Das aber ist offenbar unmöglich, da beide Functionen für $t = 0$ und $t = 1$ ungerade Zahlen sind.
